\section*{الفصل \ref{ch:b2:planet-formation-meteorites} --- تكوّن الكواكب والنيازك}

\begin{solution}{exo:b2:planet-formation-meteorites:1}
الغاز: \qty{2200}{g.cm^{-2}} عند \qty{1}{au}؛ وعند \qty{5.2}{au}،
$2200\times5.2^{-3/2} = 2200/11.86 = \qty{186}{g.cm^{-2}}$. المواد الصلبة:
$33\times0.33 = \qty{10.9}{g.cm^{-2}}$ عند \qty{1}{au}؛
$33/11.86 = \qty{2.78}{g.cm^{-2}}$ عند \qty{5.2}{au}. وعند \qty{5.2}{au} يفوق الغاز
المواد الصلبة كتلةً بعامل $186/2.78 = 67$.
\end{solution}

\begin{solution}{exo:b2:planet-formation-meteorites:2}
الألومنيوم (\qty{1653}{K}، فائق العصيان)، المغنيسيوم (\qty{1336}{K}، شائع)،
الحديد (\qty{1334}{K}، شائع)، الصوديوم (\qty{958}{K}، متوسط التطاير)، الكبريت
(\qty{664}{K}، متطاير)، الأكسجين في صورة جليد الماء (\qty{180}{K}، شديد التطاير).
والدخائل CAI، وهي أولى المتكاثفات، غنية بالألومنيوم (وبالكالسيوم والتيتانيوم)؛
والأرض فقيرة بالعناصر المتطايرة، الصوديوم والكبريت والماء، نسبةً إلى
المغنيسيوم والألومنيوم.
\end{solution}

\begin{solution}{exo:b2:planet-formation-meteorites:3}
السقطات: الكوندريتات العادية $856/1107 = \qty{77.3}{\percent}$، \omterm{def:b2:planet-formation-meteorites:classes}{والنيازك الحديدية}
$49/1107 = \qty{4.4}{\percent}$. اللقيات: $38647/44534 = \qty{86.8}{\percent}$
و$1020/44534 = \qty{2.3}{\percent}$. لقيات القارة القطبية الجنوبية: $22096/44534 =
\qty{49.6}{\percent}$. السقطات: فكل سقطة مرصودة تُجمع أيًّا كان
صنفها، أمّا اللقية فتتوقّف على مدى صمود \omterm{def:b2:planet-formation-meteorites:fall-find}{النيزك} على الأرض،
وعلى مظهره وعلى الأماكن التي يبحث فيها الناس، ونصف اللقيات يأتي من نوع واحد من
الأماكن.
\end{solution}

\begin{solution}{exo:b2:planet-formation-meteorites:4}
$\Delta t = 1.035\ln(\num{5.23e-5}/\num{1.0e-5}) = 1.035\ln 5.23 =
1.035\times1.654 = \qty{1.71}{Myr}$ بعد الدخائل CAI.
\end{solution}

\begin{solution}{exo:b2:planet-formation-meteorites:5}
$4\pi\Sigma_1 r_1^2 = 4\pi\times\num{2.2e4}\times(\num{1.496e11})^2 =
\qty{6.19e27}{kg}$. داخل \qty{30}{au}: $\times\sqrt{30} = 5.48$، ومنه
\qty{3.39e28}{kg} $= 0.017$ كتلة شمسية؛ وداخل \qty{100}{au}: $\times10$،
\qty{6.19e28}{kg} $= 0.031$ كتلة شمسية. وكلتاهما داخل المدى المرصود،
وأقرب إلى حدّه الأعلى: فقد كان السديم قرصًا عاديًا.
\end{solution}

\begin{solution}{exo:b2:planet-formation-meteorites:6}
$r = (120/T)^{7/3}$ بالوحدة الفلكية: للماء، $(120/182)^{7/3} = 0.659^{7/3} =
\qty{0.38}{au}$؛ وللميثان، $(120/41)^{7/3} = 2.93^{7/3} = \qty{12.3}{au}$.
والحبيبة المعرّضة لضوء الشمس المباشر هي الطبقة السطحية للقرص؛ أمّا المستوى الأوسط،
المحجوب بالغبار الذي فوقه، فأبرد، فيبقى الجليد أقرب إلى النجم. وقد تكوّنت
\omterm{def:b2:planet-formation-meteorites:planetesimal}{الكويكبات البدائية} في المستوى الأوسط، فينطبق تقدير المستوى الأوسط، بشرط أن يكون
القرص سلبيًا؛ والتسخين بالتراكم يزيح \omterm{def:b2:planet-formation-meteorites:snow-line}{خط الثلج} نحو الخارج.
\end{solution}

\begin{solution}{exo:b2:planet-formation-meteorites:7}
عند \qty{1}{au}: $r^2 = \num{2.238e22}$، $\Sigma_p = \qty{109}{kg.m^{-2}}$، ومنه
$2\pi b r^2\Sigma_p = 31.42\times\num{2.238e22}\times109 = \num{7.66e25}$؛
$(\num{7.66e25})^{3/2} = \num{6.71e38}$، مقسومًا على
$(3\times\num{1.989e30})^{1/2} = \num{2.443e15}$: \qty{2.75e23}{kg}
$= 0.046$ كتلة أرضية. وعند \qty{5.2}{au}: $r^2 = \num{6.05e23}$،
$\Sigma_p = \qty{27.8}{kg.m^{-2}}$، $2\pi b r^2\Sigma_p = \num{5.29e26}$،
$M = \num{1.22e40}/\num{2.443e15} = \qty{4.98e24}{kg} = 0.83$ كتلة أرضية.
ومع مواد صلبة أكثر بعامل \num{3.5} تنمو $M$ بعامل $3.5^{3/2} = 6.55$: أي 5.5 كتل
أرضية.
\end{solution}

\begin{solution}{exo:b2:planet-formation-meteorites:8}
لكل كيلوغرام $\num{1.01e19}$ ذرّة، ومنه
$\num{1.01e19}\times\num{5.00e-13} = \qty{5.05e6}{J.kg^{-1}}$. ويكون الارتفاع في درجة
الحرارة $\num{5.05e6}/837 = \qty{6030}{K}$. ويتطلّب التسخين من \num{250} إلى
\qty{1700}{K} مقدار $837\times1450 = \qty{1.21e6}{J.kg^{-1}}$: فقد وفّر الألومنيوم 26
نحو أربعة أضعاف ذلك.
\end{solution}

\begin{solution}{exo:b2:planet-formation-meteorites:9}
$\num{5.07e10}/\num{7.49e7} = 677$. ومتوسط عمر الحديد 60 هو
$2.62/\ln2 = \qty{3.78}{Myr}$. والقدرتان الابتدائيتان لكل غرام:
$\num{5.07e10}/1.035 = \num{4.90e10}$ و$\num{7.49e7}/3.78 = \num{1.98e7}$ إرغ لكل
غرام في المليون سنة، أي نسبة قدرها نحو \num{2470}: فقد ساد الألومنيوم 26 بفارق كبير.
\end{solution}

\begin{solution}{exo:b2:planet-formation-meteorites:10}
تتساوى القدرتان حين $\num{4.90e10}\,e^{-t/1.035} = \num{1.98e7}\,e^{-t/3.78}$،
أي $t(1/1.035 - 1/3.78) = \ln(\num{4.90e10}/\num{1.98e7}) = \ln 2475 = 7.81$؛
ومع $0.966 - 0.265 = 0.702$ لكل مليون سنة، $t = \qty{11.1}{Myr}$. وحينئذ
لا يبقى للحديد 60 إلا $\num{7.49e7}\,e^{-11.1/3.78} = \num{7.49e7}\times0.053 =
\qty{4.0e6}{erg.g^{-1}}$، أي \qty{400}{J.kg^{-1}}: وهذا أقل بآلاف المرات مما
يلزم لصهر أي شيء.
\end{solution}

\begin{solution}{exo:b2:planet-formation-meteorites:11}
$\tau = 1.035\times\num{3.156e13} = \qty{3.267e13}{s}$
و$H_0\tau = \qty{6.21e6}{J.kg^{-1}}$. من أجل \qty{1200}{K}:
$t_a^* = 1.035\ln\bigl(\num{6.21e6}/(837\times950)\bigr) = 1.035\ln7.81 =
\qty{2.13}{Myr}$؛ ومن أجل \qty{1700}{K}: $1.035\ln\bigl(\num{6.21e6}/(837\times1450)\bigr)
= 1.035\ln5.12 = \qty{1.69}{Myr}$. وخفض $H_0$ إلى النصف يطرح $\tau\ln2 =
\qty{0.72}{Myr}$ من كليهما، لأن $t_a^* = \tau\ln H_0 + \tau\ln(\tau/c(T_c -
T_0))$ ولا يتغيّر إلا الحدّ الأول: فالانزياح عمر نصف واحد، أيًّا كانت
درجة الحرارة الحرجة.
\end{solution}

\begin{solution}{exo:b2:planet-formation-meteorites:12}
$\kappa = 2.1/(3300\times837) = \qty{7.60e-7}{m^2.s^{-1}}$
و$\sqrt{\kappa\tau} = \sqrt{\num{7.60e-7}\times\num{3.267e13}} = \qty{4980}{m}$،
أي نحو \qty{5}{km}. ومركز جسم قطره \qty{100}{km} على عمق \qty{50}{km}،
أي عشرة أضعاف هذه المسافة: فلا تستطيع حرارته الإفلات ما دام الألومنيوم 26 يتحلّل،
وتنطبق القضية. أمّا الجسم الذي قطره \qty{2}{km} فيفقد حرارته بالسرعة التي
ينتجها بها. ويتطلّب الانصهار نصف قطر يبلغ بضعة أضعاف \qty{5}{km}، أي قطرًا من
نحو 10 إلى \qty{20}{km} أو أكثر، كما تجد النماذج.
\end{solution}

\begin{solution}{pb:b2:planet-formation-meteorites:1}
\textbf{1.} سقطة: فقد شوهد وصوله والتُقط في الصباح التالي.
فهو لم يتعرّض للتجوية، ويُضاف إلى العيّنة غير المنحازة لما يصل إلى الأرض.

\textbf{2.} $856/1107 = 0.77$: نحو ثلاث فرص من أربع.

\textbf{3.} كوندرولات قطرها \qty{0.3}{mm}، و\qty{9}{\percent} من الفلز، وأوليفين فيه
\qty{18}{\percent} من الفاياليت: كوندريت H (\qty{10}{\percent} من الفلز، Fa من 16 إلى 20).

\textbf{4.} \omterm{def:b2:planet-formation-meteorites:chondrule}{الكوندرولات} الباهتة والمعادن المتجانسة تعني تحوّلًا حراريًا: النمطان
5 و6. وقد بلغت الكوندريتات العادية على الأكثر نحو \qty{1220}{K}، أي
\qty{947}{\degreeCelsius}.

\textbf{5.} لا تحوي الكوندريتات العادية إلا من 0.1 إلى \qty{1}{\percent} من الدخائل CAI
حجمًا. والحجر مكوّن من كوندرولات وأرضية تجمّعت بعد الدخائل CAI، التي
هي أقدم المواد الصلبة: فكل صخر يحويها أحدث منها.

\textbf{6.} تلقّى كل معدن عند تكوّنه المغنيسيوم نفسه، ومعه ألومنيوم 26
تناسبًا مع ألومنيومه؛ وقد تحوّل الألومنيوم 26 منذ ذلك الحين إلى
مغنيسيوم 26، فكلما كان في المعدن ألومنيوم أكثر لكل مغنيسيوم، اكتسب
مغنيسيوم 26 أكثر.

\textbf{7.} $(0.139590 - 0.139500)/(15.0 - 0.05) = \num{0.000090}/14.95 =
\num{6.02e-6}$: وهو النسبة $\ce{^{26}Al}/\ce{^{27}Al}$ في \omterm{def:b2:planet-formation-meteorites:chondrule}{الكوندرول} حين
تكوّنه.

\textbf{8.} $0.139500 + \num{6.02e-6}\times3.00 = 0.139518$: على المستقيم.

\textbf{9.} $\Delta t = 1.035\ln(\num{5.23e-5}/\num{6.02e-6}) = 1.035\ln8.69 =
1.035\times2.162 = \qty{2.24}{Myr}$.

\textbf{10.} $4567.5 - 2.2 = \qty{4565.3}{\Ma}$.

\textbf{11.} $\num{5.23e-5}\,e^{-10/1.035} = \num{5.23e-5}\times\num{6.4e-5} =
\num{3.3e-9}$، وهذا أقل بكثير من أن يُقاس فائضه من المغنيسيوم 26. أمّا الهافنيوم 182 (عمر النصف
\qty{8.9}{Myr}) أو المنغنيز 53 (\qty{3.7}{Myr}) فيناسبان أحداثًا وقعت بعد نحو عشرة ملايين سنة
من الدخائل CAI.

\textbf{12.} $\kappa = 2.1/(3300\times837) = \qty{7.60e-7}{m^2.s^{-1}}$؛
$R^2/\kappa = \num{4e8}/\num{7.60e-7} = \qty{5.26e14}{s} = \qty{16.7}{Myr}$، أي ستة عشر
ضعفًا لمتوسط عمر الألومنيوم 26: فالمركز يحتفظ بحرارته.

\textbf{13.} $\num{1.9e-7}\,e^{-2.4/1.035} = \num{1.9e-7}\times0.0984 =
\qty{1.87e-8}{W.kg^{-1}}$.

\textbf{14.} $\num{1.87e-8}\times\num{3.267e13} = \qty{6.11e5}{J.kg^{-1}}$.

\textbf{15.} $\Delta T = \num{6.11e5}/837 = \qty{730}{K}$؛ والذروة
$250 + 730 = \qty{980}{K}$.

\textbf{16.} فوق \qty{750}{K}، فقد تلبّد الصخر وتحوّل؛ وتحت
\qty{1200}{K}، فلم ينصهر الفلز ولا الكبريتيد. فالحجر كوندريت مسخَّن،
من نمط صخري مرتفع، دون \qty{1220}{K} الواردة في السؤال 4: وهذا متّسق.

\textbf{17.} $1.035\ln\bigl(\num{6.21e6}/(837\times950)\bigr) = \qty{2.13}{Myr}$:
أي أبكر بنحو \qty{0.3}{Myr}.

\textbf{18.} الارتفاع هو $\Delta T(1 - e^{-(t - t_a)/\tau})$؛ وبلوغ \qty{90}{\percent}
يتطلّب $e^{-(t - t_a)/\tau} = 0.1$، أي $t - t_a = 1.035\ln10 = \qty{2.38}{Myr}$، أي
بعد \qty{4.8}{Myr} من الدخائل CAI.

\textbf{19.} $\Delta T = 7420\,e^{-0.8/1.035} = 7420\times0.462 = \qty{3430}{K}$،
والذروة \qty{3680}{K}، أعلى بكثير من \qty{1700}{K} التي تكون عندها السيليكات منصهرة.
فقد انصهر الجسم (وأبقت حرارةُ الانصهار درجةَ حرارته منخفضة)
وتمايز إلى نواة فلزية ووشاح وقشرة سيليكاتيين.

\textbf{20.} أكثر بعامل $e^{(2.4 - 0.8)/1.035} = e^{1.546} = 4.7$.

\textbf{21.} يتكاثف الزنك والرصاص قرب \qty{727}{K}، فهما متوسطا التطاير؛ أمّا المغنيسيوم
(\qty{1336}{K}) والألومنيوم (\qty{1653}{K}) فيتكاثفان أبكر بكثير. وقد تجمّعت المواد الصلبة
في منطقة الحجر قبل أن يبرد الغاز بما يكفي لتكاثف كل الزنك
والرصاص.

\textbf{22.} $33\times0.33\times2.5^{-3/2} = 10.89/3.95 = \qty{2.76}{g.cm^{-2}}$،
أي \qty{27.6}{kg.m^{-2}}.

\textbf{23.} $r^2 = (3.74\times10^{11})^2 = \num{1.40e23}$؛
$2\pi\times5\times\num{1.40e23}\times27.6 = \num{1.21e26}$؛
$(\num{1.21e26})^{3/2}/\num{2.443e15} = \qty{5.45e23}{kg}$، أي $0.091$ كتلة أرضية.
أمّا الجسم الأم، $\tfrac43\pi(\num{2e4})^3\times3300 = \qty{1.1e17}{kg}$، فأخفّ بخمسة
ملايين مرة: إنه \omterm{def:b2:planet-formation-meteorites:planetesimal}{كويكب بدائي}، لا جنين.

\textbf{24.} بلغت ذروة مركز الجسم الأم نحو
\textbf{\qty[text-series-to-math=true]{707}{\degreeCelsius}} ($980 - 273$): وهي حرارة كافية
لإعادة تبلّر \omterm{def:b2:planet-formation-meteorites:chondrule}{الكوندريت}، وأبرد من أن تصهر فلزه.
\end{solution}
